\documentclass[../../main.tex]{subfiles}
\begin{document}
\section{习题}\label{section:ch09-习题}

\begin{exercise}\label{exercise:ch09-01}
假设 $M$ 是光滑流形，$X\in\mathfrak{X}(M)$，$\gamma$ 是 $X$ 的极大积分曲线。
\begin{enumerate}[(a)]
\item 我们说 $\gamma$ 是\textbf{周期的}，若存在数 $T>0$ 使得对所有 $t\in\mathbb{R}$，$\gamma(t+T)=\gamma(t)$。证明下列恰有一个成立：
\begin{itemize}
\item $\gamma$ 是常值的。
\item $\gamma$ 是单射的。
\item $\gamma$ 是周期且非常值的。
\end{itemize}
\item 证明若 $\gamma$ 周期且非常值，则存在唯一正数 $T$（称为 $\gamma$ 的\textbf{周期}），使得 $\gamma(t)=\gamma(t_0)$ 当且仅当对某 $k\in\mathbb{Z}$，$t-t_0=kT$。
\item 证明 $\gamma$ 的像是 $M$ 的浸入子流形，微分同胚于 $\mathbb{R}$、$S^1$ 或 $\mathbb{R}^0$。
\end{enumerate}
（用于第 398、560 页。）
\end{exercise}

\begin{exercise}\label{exercise:ch09-02}
假设 $M$ 是光滑流形，$S\subseteq M$ 是浸入子流形，$V$ 是 $M$ 上与 $S$ 相切的光滑向量场。
\begin{enumerate}[(a)]
\item 证明对任意满足 $\gamma(t_0)\in S$ 的 $V$ 的积分曲线 $\gamma$，存在 $\varepsilon>0$ 使得 $\gamma((t_0-\varepsilon;t_0+\varepsilon))\subseteq S$。
\item 现假设 $S$ 正常嵌入。证明每条与 $S$ 相交的积分曲线都包含于 $S$。
\item 若 $S$ 不闭，给出 (b) 的反例。
\end{enumerate}
（用于第 243、491 页。）
\end{exercise}

\begin{exercise}\label{exercise:ch09-03}
计算 $\mathbb{R}^2$ 上下列每个向量场的流：
\begin{enumerate}[(a)]
\item $V=y\dfrac{\partial}{\partial x}+\dfrac{\partial}{\partial y}$。
\item $W=x\dfrac{\partial}{\partial x}+2y\dfrac{\partial}{\partial y}$。
\item $X=x\dfrac{\partial}{\partial x}-y\dfrac{\partial}{\partial y}$。
\item $Y=x\dfrac{\partial}{\partial y}+y\dfrac{\partial}{\partial x}$。
\end{enumerate}
\end{exercise}

\begin{exercise}\label{exercise:ch09-04}
对任意整数 $n\geqslant 1$，在奇维球面 $S^{2n-1}\subseteq\mathbb{C}^n$ 上由 $\theta(t;z)=e^{it}z$ 定义流。证明 $\theta$ 的无穷小生成元是 $S^{2n-1}$ 上光滑非零向量场。[注：在 $n=2$ 的情形，$X$ 的积分曲线是习题 3-6 的曲线 $\gamma_z$，故这提供了每条 $\gamma_z$ 光滑的更简单证明。]（用于第 435 页。）
\end{exercise}

\begin{exercise}\label{exercise:ch09-05}
假设 $M$ 是光滑紧流形，允许处处非零的光滑向量场。证明存在光滑映射 $F\colon M\to M$，同伦于恒等且无不动点。
\end{exercise}

\begin{exercise}\label{exercise:ch09-06}
证明\cref{lemma:ch09-19}（逃逸引理）。
\end{exercise}

\begin{exercise}\label{exercise:ch09-07}
设 $M$ 是连通光滑流形。证明 $M$ 的微分同胚群在 $M$ 上传递作用：即对任意 $p,q\in M$，存在微分同胚 $F\colon M\to M$ 使 $F(p)=q$。[提示：先证明若 $p,q\in B^n$（$\mathbb{R}^n$ 中的开单位球），则 $B^n$ 上存在紧支集光滑向量场，其流 $\theta$ 满足 $\theta_1(p)=q$。]
\end{exercise}

\begin{exercise}\label{exercise:ch09-08}
设 $M$ 是光滑流形，$S\subseteq M$ 是紧嵌入子流形。假设 $V\in\mathfrak{X}(M)$ 是处处不与 $S$ 相切的光滑向量场。证明存在 $\varepsilon>0$ 使得 $V$ 的流限制为光滑嵌入 $\Phi\colon(-\varepsilon,\varepsilon)\times S\to M$。
\end{exercise}

\begin{exercise}\label{exercise:ch09-09}
假设 $M$ 是光滑流形，$S\subseteq M$ 是嵌入超曲面（不一定紧）。进一步假设存在定义在 $S$ 邻域上、处处不与 $S$ 相切的光滑向量场 $V$。证明 $S$ 在 $M$ 中有微分同胚于 $(-1;1)\times S$ 的邻域，且该微分同胚限制为显然的等同 $\{0\}\times S\leftrightarrow S$。[提示：用流出定理的记号，证明 $\mathcal{O}_\delta\leftrightarrow\mathcal{O}_1$。]
\end{exercise}

\begin{exercise}\label{exercise:ch09-10}
对习题 9-3 中的每个向量场，在 $(1;0)$ 邻域中寻找光滑坐标，使得给定向量场是坐标向量场。
\end{exercise}

\begin{exercise}\label{exercise:ch09-11}
证明\cref{theorem:ch09-24}（边界流出定理）。[提示：先在边界坐标中定义 $\Phi$，用唯一性将局部定义粘接起来。为得到嵌入，确保 $\delta(p)$ 不超过从 $p$ 出发的积分曲线首次击中边界（若有的话）的时间的一半。]
\end{exercise}

\begin{exercise}\label{exercise:ch09-12}
假设 $M_1$ 和 $M_2$ 是连通光滑 $n$-流形，$M_1\# M_2$ 是它们的光滑连通和（见\cref{example:ch09-31}）。证明 $M_1\# M_2$ 上的光滑结构可选得使存在开子集 $\widehat{M}_1,\widehat{M}_2\subseteq M_1\# M_2$，分别微分同胚于 $M_1\smallsetminus\{p_1\}$ 和 $M_2\smallsetminus\{p_2\}$，满足 $\widehat{M}_1\cup\widehat{M}_2=M_1\# M_2$ 且 $\widehat{M}_1\cap\widehat{M}_2$ 微分同胚于 $(-1;1)\times S^{n-1}$。（用于第 465 页。）
\end{exercise}

\begin{exercise}\label{exercise:ch09-13}
证明若允许 $M$ 为带边光滑流形，\cref{theorem:ch05-29}和\cref{theorem:ch05-53}(b)（将光滑映射的 codomain 限制到子流形或带边子流形）的结论仍然成立。
\end{exercise}

\begin{exercise}\label{exercise:ch09-14}
用二重证明\cref{theorem:ch06-18}（Whitney 浸入定理）在 $M$ 有非空边界的情形。
\end{exercise}

\begin{exercise}\label{exercise:ch09-15}
证明\cref{theorem:ch09-35}（边界正则点附近的标准型）。[提示：将 $M$ 视为其二重中的正则区域，从 $x^n$ 是 $D(M)$ 中 $\partial M$ 的局部定义函数的坐标开始。]
\end{exercise}

\begin{exercise}\label{exercise:ch09-16}
给出 $\mathbb{R}^2$ 上光滑向量场 $V,\widetilde{V},W$ 的例子，使得沿 $x$ 轴 $V=\widetilde{V}=\partial/\partial x$，但在原点处 $L_VW\neq L_{\widetilde{V}}W$。[注：这表明要计算 $(L_VW)_p$ 确实需要知道向量场 $V$；仅知道向量 $V_p$，甚至仅知道 $V$ 沿 $V$ 的某条积分曲线的值，都是不够的。]
\end{exercise}

\begin{exercise}\label{exercise:ch09-17}
对下列 $\mathbb{R}^3$ 上的 $k$ 元向量场，要么在 $(1;0;0)$ 邻域中寻找光滑坐标 $(s^1;s^2;s^3)$ 使得对 $i=1,\dots,k$，$V_i=\partial/\partial s^i$，要么解释为何不存在。
\begin{enumerate}[(a)]
\item $k=2$：$V_1=\dfrac{\partial}{\partial x}$；$V_2=\dfrac{\partial}{\partial x}+\dfrac{\partial}{\partial y}$。
\item $k=2$：$V_1=(x+1)\dfrac{\partial}{\partial x}-(y+1)\dfrac{\partial}{\partial y}$；$V_2=(x+1)\dfrac{\partial}{\partial x}+(y+1)\dfrac{\partial}{\partial y}$。
\item $k=3$：$V_1=x\dfrac{\partial}{\partial y}-y\dfrac{\partial}{\partial x}$；$V_2=y\dfrac{\partial}{\partial z}-z\dfrac{\partial}{\partial y}$；$V_3=z\dfrac{\partial}{\partial x}-x\dfrac{\partial}{\partial z}$。
\end{enumerate}
\end{exercise}

\begin{exercise}\label{exercise:ch09-18}
在平面上定义向量场 $X$ 和 $Y$ 为
\begin{align*}
X=x\frac{\partial}{\partial x}-y\frac{\partial}{\partial y},\qquad Y=x\frac{\partial}{\partial y}+y\frac{\partial}{\partial x}.
\end{align*}
计算 $X$ 和 $Y$ 的流 $\theta,\psi$，并通过寻找显式的包含 $0$ 的开区间 $J$ 和 $K$ 来验证流不交换，使得 $\theta_s\circ\psi_t$ 和 $\psi_t\circ\theta_s$ 对所有 $(s;t)\in J\times K$ 都有定义，但对某些此类 $(s;t)$ 它们不相等。
\end{exercise}

\begin{exercise}\label{exercise:ch09-19}
设 $M$ 为去掉 $z$ 轴的 $\mathbb{R}^3$。定义 $V,W\in\mathfrak{X}(M)$ 为
\begin{align*}
V=\frac{\partial}{\partial x}-\frac{y}{x^2+y^2}\frac{\partial}{\partial z},\qquad W=\frac{\partial}{\partial y}+\frac{x}{x^2+y^2}\frac{\partial}{\partial z},
\end{align*}
设 $\theta$ 和 $\psi$ 分别为 $V$ 和 $W$ 的流。证明 $V$ 和 $W$ 交换，但存在 $p\in M$ 和 $s,t\in\mathbb{R}$ 使得 $\theta_t\circ\psi_s(p)$ 和 $\psi_s\circ\theta_t(p)$ 都有定义但不相等。
\end{exercise}

\begin{exercise}\label{exercise:ch09-20}
假设 $M$ 是紧光滑流形，$V\colon J\times M\to TM$ 是 $M$ 上的光滑含时向量场。证明 $V$ 的含时流的定义域是整个 $J\times J\times M$。
\end{exercise}

\begin{exercise}\label{exercise:ch09-21}
设 $M$ 是光滑流形。$M$ 的\textbf{光滑同痕}是光滑映射 $H\colon M\times J\to M$，其中 $J\subseteq\mathbb{R}$ 是区间，使得对每个 $t\in J$，由 $H_t(p)=H(p;t)$ 定义的映射 $H_t\colon M\to M$ 是微分同胚。（特别地，若 $J$ 是单位区间，则 $H$ 是经过微分同胚从 $H_0$ 到 $H_1$ 的同伦。）本问题表明光滑同痕与含时流密切相关。
\begin{enumerate}[(a)]
\item 假设 $J\subseteq\mathbb{R}$ 是开区间，$H\colon M\times J\to M$ 是光滑同痕。证明由 $V(t;p)=\frac{\partial}{\partial t}H(p;t)$ 定义的映射 $V\colon J\times M\to TM$ 是 $M$ 上的光滑含时向量场，其含时流由 $\Psi(t;t_0;p)=H_t\circ H_{t_0}^{-1}(p)$ 给出，定义域 $J\times J\times M$。
\item 反过来，假设 $J$ 是开区间，$V\colon J\times M\to M$ 是 $M$ 上的光滑含时向量场，其含时流定义在 $J\times J\times M$ 上。对任意 $t_0\in J$，证明由 $H(t;p)=\Psi(t;t_0;p)$ 定义的映射 $H\colon M\times J\to M$ 是 $M$ 的光滑同痕。
\end{enumerate}
\end{exercise}

\begin{exercise}\label{exercise:ch09-22}
下面是 $\mathbb{R}^2$ 中的三个 Cauchy 问题。对每个，求初始子流形邻域中的显式解 $u(x;y)$：
\begin{enumerate}[(a)]
\item $y\dfrac{\partial u}{\partial x}+\dfrac{\partial u}{\partial y}=x$；$u(x;0)=\sin x$。
\item $\dfrac{\partial u}{\partial x}+y\dfrac{\partial u}{\partial y}=0$；$u(x;1)=e^{-x}$。
\item $\dfrac{\partial u}{\partial x}+u\dfrac{\partial u}{\partial y}=y$；$u(0;y)=0$。
\end{enumerate}
\end{exercise}

\begin{exercise}\label{exercise:ch09-23}
再次考虑习题 9-22 中的 Cauchy 问题。证明 (a) 有唯一全局解；(b) 有全局解但不唯一；(c) 无全局解。[提示：考虑哪些特征曲线与初始子流形相交、哪些不相交。]
\end{exercise}

\begin{exercise}\label{exercise:ch09-24}
证明 Euler 齐次函数定理（习题 8-2）的逆：若 $f\in C^\infty(\mathbb{R}^n\smallsetminus\{0\})$ 满足 $Vf=cf$，其中 $V$ 是 Euler 向量场，$c\in\mathbb{R}$，则 $f$ 是 $c$ 次正齐次函数。
\end{exercise}

\end{document}